% !TEX program = xelatex
\documentclass[11pt,a4paper]{article}
\usepackage{algo-preamble}

\title{\textbf{L03 \textperiodcentered{} Διαίρει \& Κυρίευε I}\\[2pt]
\large Mergesort \& Master Theorem}
\author{Algorithms Class Hub}
\date{\small Αλγόριθμοι και Πολυπλοκότητα (K17) \textperiodcentered{} ΕΚΠΑ}

\begin{document}
\maketitle

\begin{center}\itshape\small
Το πρώτο σχήμα σχεδίασης αλγορίθμων: διαίρει και κυρίευε. Bubble/insertion sort,
mergesort, τέσσερις διαφορετικές αποδείξεις του $O(n \log n)$, πύργος του Hanoi,
και το Master Theorem.
\end{center}

\section{Τι θα δούμε}

Αυτή είναι η \textbf{πρώτη} τεχνική σχεδίασης αλγορίθμων του μαθήματος. Στόχοι:

\begin{enumerate}
  \item Να καταλάβουμε τι σημαίνει \textbf{«διαίρει και κυρίευε»} ως σχήμα --- όχι
  ως ένας συγκεκριμένος αλγόριθμος, αλλά ως οικογένεια αλγορίθμων.
  \item Να δούμε το πιο γνωστό παράδειγμα: τη \textbf{συγχωνευτική ταξινόμηση
  (mergesort)}, να αποδείξουμε ότι τρέχει σε $O(n \log n)$ με \textbf{τέσσερις
  διαφορετικούς τρόπους}, και να καταλάβουμε γιατί κάθε τρόπος είναι χρήσιμος.
  \item Να βγάλουμε ένα γενικό εργαλείο --- το \textbf{Master Theorem} --- που
  λύνει ολόκληρη οικογένεια αναδρομικών σχέσεων χωρίς απόδειξη από την αρχή.
  \item Να δούμε ότι D\&C δεν είναι «μόνο για ταξινόμηση» --- εφαρμόζεται και σε
  \textbf{πύργο του Hanoi}, και σε πολλά άλλα προβλήματα που θα δούμε στα επόμενα
  μαθήματα.
\end{enumerate}

\section{Σύντομη υπενθύμιση από το δεύτερο μάθημα}

Ας συμφωνήσουμε στη γλώσσα πριν προχωρήσουμε. Από το δεύτερο μάθημα:

\[
O(f(n)) = \{T : \exists n_0,\ \exists c > 0,\ \forall n \ge n_0,\ T(n) \le c
\cdot f(n)\}
\]
\[
\Omega(f(n)) = \{T : \exists n_0,\ \exists c > 0,\ \forall n \ge n_0,\ T(n) \ge c
\cdot f(n)\}
\]
\[
\Theta(f(n)) = O(f(n)) \cap \Omega(f(n))
\]

Και εναλλακτικά μέσω ορίων:

\begin{center}
\begin{tabular}{@{}l l@{}}
\toprule
$\lim_{n \to \infty} f(n)/g(n)$ & Συμβολισμός \\
\midrule
$\alpha \ne 0$ & $f \in \Theta(g)$ \\
$0$            & $f \in O(g)$, $f \notin \Theta(g)$, δηλαδή $f \in o(g)$ \\
$\infty$       & $g \in O(f)$, $g \notin \Theta(f)$, δηλαδή $f \in \omega(g)$ \\
\bottomrule
\end{tabular}
\end{center}

Και η \textbf{κλίμακα πολυπλοκοτήτων}:
\[
1 \prec \log n \prec n \prec n \log n \prec n^2 \prec n^3 \prec \dots \prec n^p
\prec 2^n \prec 3^n \prec n!
\]
Η \textbf{πολυωνυμική} πολυπλοκότητα είναι \textbf{καλή}· η \textbf{εκθετική}
είναι \textbf{κακή}.

\begin{keypoint}
\textbf{Βέλτιστος αλγόριθμος} (επανάληψη): αν για ένα πρόβλημα $\Pi$ υπάρχει
\textbf{κάτω φράγμα} $\Omega(f(n))$ --- δηλαδή αποδεικνύεται ότι κάθε αλγόριθμος
που λύνει το $\Pi$ απαιτεί τουλάχιστον $\Omega(f(n))$ βήματα --- και έχουμε
αλγόριθμο σε $O(f(n))$, τότε ο αλγόριθμός μας είναι \textbf{ασυμπτωτικά
βέλτιστος}. Στο τέλος αυτής της διάλεξης θα ξέρουμε ότι η mergesort είναι ακριβώς
αυτή η περίπτωση.
\end{keypoint}

\section{Ζέσταμα: μέγιστο στοιχείο σε πίνακα}

Πριν πάμε στο σύνθετο πρόβλημα της ταξινόμησης, ας δούμε ένα πολύ απλό για να
εξασκηθούμε με την ανάλυση. \textbf{Πρόβλημα:} δοθέντος πίνακα
$a_1, a_2, \dots, a_n$, να βρεθεί το μέγιστο στοιχείο.

\begin{examplebox}
\textbf{Αλγόριθμος 1 --- «κράτα τον τρέχοντα μέγιστο».}
\begin{codebox}
\begin{verbatim}
v <- a_1
Για i = 2 έως n:
  Εάν v < a_i: v <- a_i
επίστρεψε v
\end{verbatim}
\end{codebox}
Στοιχειώδης πράξη: σύγκριση. Κάνουμε \textbf{ακριβώς $n-1$ συγκρίσεις},
ανεξάρτητα από την είσοδο. Άρα $\Theta(n)$.

\textbf{Αλγόριθμος 2 --- «ωθώ τον μέγιστο στο τέλος».}
\begin{codebox}
\begin{verbatim}
Για i = 1 έως n-1:
  Εάν a_i > a_{i+1}: ενάλλαξε
επίστρεψε a_n
\end{verbatim}
\end{codebox}
Πάλι $n-1$ συγκρίσεις. $\Theta(n)$.

\textbf{Παρατήρηση.} Και οι δύο αλγόριθμοι έχουν την ίδια ασυμπτωτική
πολυπλοκότητα, αλλά ο 1ος είναι σαφώς προτιμότερος: ο 2ος \textit{τροποποιεί} τον
πίνακα εισόδου, ενώ ο 1ος όχι. Όταν δύο αλγόριθμοι έχουν το ίδιο ασυμπτωτικό
κόστος, ψάχνουμε \textbf{άλλα κριτήρια} (μνήμη, ευκολία υλοποίησης, σταθερότητα,
παράλληλη εκτέλεση).
\end{examplebox}

\section{Το πραγματικό πρόβλημα: ταξινόμηση}

\begin{quoteblock}
\textbf{Πρόβλημα.} Δοθέντος πίνακα από στοιχεία, να τα ταξινομήσουμε (με βάση
κάποιο κριτήριο σύγκρισης).
\end{quoteblock}

\subsection{Προφανείς εφαρμογές}

\begin{itemize}
  \item Οργάνωση μιας \textbf{playlist} ή αλφαβητική λίστα τηλεφώνων.
  \item Απαρίθμηση αρχείων σε έναν φάκελο αντίστροφα από την ημερομηνία.
  \item κ.ο.κ.
\end{itemize}

\subsection{Λιγότερο προφανείς εφαρμογές}

Μια \textbf{ταξινομημένη} λίστα είναι το «κλειδί» για πολλά:
\begin{itemize}
  \item \textbf{Εύρεση μέσου στοιχείου.} Σε ταξινομημένη λίστα, το μεσαίο
  στοιχείο είναι σε $O(1)$.
  \item \textbf{Εύρεση κοντινού ζευγαριού.} Σε ταξινομημένη λίστα, το πιο κοντινό
  ζευγάρι αριθμών είναι κάποια διπλανή θέση --- $O(n)$.
  \item \textbf{Δυαδική αναζήτηση.} Όπως είδαμε στο πρώτο μάθημα, σε ταξινομημένη
  λίστα η αναζήτηση είναι $O(\log n)$ αντί για $O(n)$.
  \item \textbf{Εντοπισμός ακραίων τιμών (outliers).} Στις άκρες της
  ταξινομημένης λίστας.
  \item \textbf{Εντοπισμός επαναλαμβανόμενων στοιχείων.} Διαδοχικά στοιχεία ή ίσα
  στοιχεία, σε $O(n)$.
\end{itemize}

Και \textbf{μη προφανείς} --- όλες αυτές οι περιοχές χρησιμοποιούν ταξινόμηση ως
υποπρόβλημα: συμπίεση δεδομένων (π.χ.\ Huffman coding, που θα δούμε αργότερα στους
άπληστους αλγορίθμους), γραφικά Η/Υ (z-buffer ταξινόμηση πολυγώνων), υπολογιστική
βιολογία (DNA alignment), και ελάχιστα συνδετικά δέντρα (αλγόριθμος Kruskal στην
ενότητα των γραφημάτων --- απαιτεί ταξινομημένες ακμές).

\begin{intuition}
Η ταξινόμηση είναι από τα \textbf{πιο μελετημένα} προβλήματα στην επιστήμη
υπολογιστών γιατί είναι \textbf{παντού}. Ένας καλός αλγόριθμος ταξινόμησης
ξεκλειδώνει βελτιώσεις σε όλα τα παραπάνω.
\end{intuition}

\section{Πρώτη προσπάθεια: bubble sort \& insertion sort}

Πριν δούμε το «έξυπνο» mergesort, ας δούμε τα κλασικά απλά.

\begin{examplebox}
\textbf{Bubble sort --- η «φυσαλίδα».}
\begin{codebox}
\begin{verbatim}
Διαδικασία bubblesort(a_1, ..., a_n):
  Για i = 1 έως n-1:
    Για j = 1 έως n-i:
      Αν a_j > a_{j+1}: ενάλλαξε a_j <-> a_{j+1}
\end{verbatim}
\end{codebox}
\textbf{Διαίσθηση.} Σε κάθε πέρασμα του εξωτερικού βρόχου, το «μεγαλύτερο»
στοιχείο «αναδύεται» (σαν φυσαλίδα) στο σωστό του τέλος. Μετά από $n-1$
περάσματα, όλα είναι στη θέση τους.

\textbf{Πολυπλοκότητα.} Διπλά φωλιασμένος βρόχος. Στο πρώτο πέρασμα κάνουμε $n-1$
συγκρίσεις, στο δεύτερο $n-2$, \dots, στο τελευταίο 1 σύγκριση. Σύνολο:
\[
(n-1) + (n-2) + \dots + 1 = \frac{n(n-1)}{2} = \Theta(n^2)
\]
Στη \textbf{χείριστη περίπτωση} (πίνακας αντίστροφα ταξινομημένος) γίνονται
\textit{όλες} οι ανταλλαγές --- $\Theta(n^2)$.
\end{examplebox}

\begin{examplebox}
\textbf{Insertion sort --- η «ενθετική».}
\begin{codebox}
\begin{verbatim}
Διαδικασία insertionsort(a_1, ..., a_n):
  Για i = 1 έως n:
    v <- a_i
    j <- i - 1
    Όσο j >= 1 και v < a_j:
      a_{j+1} <- a_j     % ολίσθηση
      j <- j - 1
    a_{j+1} <- v
\end{verbatim}
\end{codebox}
\textbf{Διαίσθηση.} Όπως ταξινομούμε τα φύλλα ενός παιχνιδιού καρτών στο χέρι:
παίρνουμε κάθε νέο φύλλο και το βάζουμε στη σωστή θέση \textit{μεταξύ} των ήδη
ταξινομημένων.

\textbf{Πολυπλοκότητα.} Στη χείριστη περίπτωση (αντίστροφα ταξινομημένος πίνακας)
πάλι $\Theta(n^2)$ --- κάθε νέο στοιχείο πρέπει να «κατέβει» σε όλο τον ήδη
ταξινομημένο πρόθεμα.

\textbf{Παρατήρηση.} Στη \textbf{βέλτιστη} περίπτωση (ήδη ταξινομημένος πίνακας)
η insertion sort είναι $\Theta(n)$ --- κάθε νέο στοιχείο μένει στη θέση του. Η
bubble sort \textit{χωρίς early-exit optimization} παραμένει $\Theta(n^2)$ ακόμα
και στην καλύτερη περίπτωση.
\end{examplebox}

\begin{warningbox}
Και οι δύο αλγόριθμοι έχουν \textbf{χείριστη πολυπλοκότητα $\Theta(n^2)$}. Για
$n = 10^6$ (ένα εκατομμύριο στοιχεία), αυτό σημαίνει $10^{12}$ πράξεις --- χρόνος
μιας \textbf{ώρας} σε σύγχρονο επεξεργαστή. Η mergesort που θα δούμε αμέσως κάνει
την ίδια δουλειά σε $20 \cdot 10^6 \approx 2 \cdot 10^7$ πράξεις --- \textbf{κάποια
δευτερόλεπτα}. Η διαφορά $n^2$ vs $n \log n$ είναι \textbf{πέντε τάξεις μεγέθους}
εκεί.
\end{warningbox}

\section{Η ιδέα: συγκρίσεις ανά δύο, μετά συγχώνευση ανά δύο}

Πριν τυποποιήσουμε, η διαίσθηση: \textbf{σπάμε} τον πίνακα σε ζευγάρια,
ταξινομούμε κάθε ζευγάρι (τετριμμένο), και μετά \textbf{συγχωνεύουμε} ζευγάρια σε
τετράδες, τετράδες σε οκτάδες, και ούτω καθεξής μέχρι να έχουμε έναν ταξινομημένο
πίνακα.

\begin{quoteblock}
\textbf{\large Divide et impera!}\\
\textit{(λατινικό: «διαίρει και βασίλευε»)}
\end{quoteblock}

\section{Διαίρει και Κυρίευε --- το γενικό σχήμα}

Η τεχνική \textbf{διαίρει και κυρίευε (D\&C)} βασίζεται σε τρία βήματα:

\begin{enumerate}
  \item \textbf{Διαιρούμε} τα δεδομένα του προβλήματός μας σε μικρότερα
  υποπροβλήματα που \textbf{δεν έχουν κοινά στοιχεία} (ξένα μεταξύ τους).
  \item \textbf{Κυριεύουμε} τα μικρότερα υποπροβλήματα \textbf{ξεχωριστά} (και
  \textbf{αναδρομικά} --- εφαρμόζοντας ξανά τον ίδιο αλγόριθμο).
  \item \textbf{Συνδυάζουμε} τις λύσεις που βρήκαμε για να λύσουμε το αρχικό
  πρόβλημα.
\end{enumerate}

\subsection{Πώς συνήθως εφαρμόζεται}

Στη συνηθισμένη μορφή:
\begin{itemize}
  \item \textbf{Διαιρούμε} δεδομένα μεγέθους $n$ σε \textbf{δύο} μικρότερα
  στιγμιότυπα μεγέθους $\lceil n/2 \rceil$ και $\lfloor n/2 \rfloor$ (ξένα μεταξύ
  τους).
  \item \textbf{Επιλύουμε} καθένα με \textbf{αναδρομική κλήση} του ίδιου
  αλγορίθμου. Θεωρούμε την αναδρομική κλήση \textbf{«μαύρο κουτί»} που μας
  επιστρέφει τη λύση.
  \item \textbf{Συνδυάζουμε} τις δύο επιμέρους λύσεις σε λύση του πλήρους
  προβλήματος.
\end{itemize}

\subsection{Γιατί δουλεύει: επαγωγική κατασκευή}

Η τεχνική D\&C βασίζεται σε \textbf{επαγωγική κατασκευή της λύσης}. Αρκεί (και
χρειάζεται) να ξέρουμε:
\begin{enumerate}
  \item \textbf{Βάση επαγωγής:} να επιλύσουμε το πρόβλημα για \textbf{μικρή
  (σταθερή) διάσταση δεδομένων} --- π.χ.\ για $n = 1$ ή $n = 2$.
  \item \textbf{Επαγωγικό βήμα:} να \textbf{συνδυάσουμε} τις ορθές λύσεις των
  ξένων υποπροβλημάτων μας σε \textbf{μία καθολική λύση} για όλα τα δεδομένα.
\end{enumerate}

\begin{keypoint}
\textbf{Η σύνδεση με μαθηματική επαγωγή είναι ακριβής.} Αν αποδείξεις (α) ότι ο
αλγόριθμος δουλεύει για $n = 1$ και (β) ότι, υποθέτοντας ότι δουλεύει για κάθε
$k < n$, δουλεύει και για $n$ --- τότε \textbf{έχεις απόδειξη ορθότητας}. Αυτό
είναι το πρότυπο απόδειξης ορθότητας για κάθε D\&C αλγόριθμο.
\end{keypoint}

\section{Συγχωνευτική ταξινόμηση --- η ιδέα}

Πώς ταξινομούμε σε αύξουσα σειρά με D\&C;

\begin{enumerate}
  \item \textbf{Βάση.} Αν ο πίνακας έχει \textbf{ένα στοιχείο}, είναι ήδη
  ταξινομημένος. Επιστρέφουμε τον ίδιο πίνακα.
  \item \textbf{Αλλιώς:} \textbf{διαιρούμε} τον πίνακα \textbf{στη μέση} (δύο ίσα
  μισά)· \textbf{ταξινομούμε} καθένα από τα δύο με μία \textbf{αναδρομική κλήση}·
  \textbf{συγχωνεύουμε} τις δύο ταξινομημένες λίστες σε μία.
\end{enumerate}

Η ιδέα είναι: «αν ξέρω να συγχωνεύσω \textbf{δύο ταξινομημένες λίστες}, και αν
ξέρω να σπάσω έναν πίνακα στη μέση, τότε μπορώ να ταξινομήσω τα πάντα» --- η
αναδρομή κάνει το υπόλοιπο.

\section{Επεξεργασμένο παράδειγμα: ταξινόμηση 11 στοιχείων}

Ακολουθούμε τον αλγόριθμο στον πίνακα $[10, 5, 8, 6, 3, 4, 7, 1, 9, 2, 11]$:

\begin{enumerate}
  \item \textbf{Είσοδος:} $[10, 5, 8, 6, 3, 4, 7, 1, 9, 2, 11]$.
  \item \textbf{Διαίρει στη μέση:} αριστερό μισό $[10, 5, 8, 6, 3, 4]$ (6
  στοιχεία), δεξί μισό $[7, 1, 9, 2, 11]$ (5 στοιχεία).
  \item \textbf{Αναδρομικά}, κάθε μισό επιστρέφει ταξινομημένο: αριστερό
  $\to [3, 4, 5, 6, 8, 10]$, δεξί $\to [1, 2, 7, 9, 11]$.
  \item \textbf{Συγχώνευση} των δύο $\to [1, 2, 3, 4, 5, 6, 7, 8, 9, 10, 11]$.
\end{enumerate}

Το σύνολο της δουλειάς ξοδεύτηκε στο βήμα 4, που είναι γραμμικό. Στην
πραγματικότητα η αναδρομή πάει βαθύτερα --- το αριστερό μισό σπάει σε δύο, καθένα
από αυτά σπάει σε δύο, \dots, μέχρι να φτάσουμε σε πίνακες μήκους 1.

\section{Η συγχώνευση δύο ταξινομημένων πινάκων}

Πριν αποδείξουμε ότι η mergesort είναι $O(n \log n)$, χρειαζόμαστε την «βασική
πράξη» που μετατρέπει τα ταξινομημένα μισά σε ταξινομημένο σύνολο: τη
\textbf{συγχώνευση}.

\begin{examplebox}
\textbf{Πρόβλημα.} Έστω $a$ και $b$ δύο \textbf{ήδη ταξινομημένοι} πίνακες με $m$
και $n$ στοιχεία αντίστοιχα. Να παράξουμε έναν ταξινομημένο πίνακα $c$ με $m+n$
στοιχεία που περιέχει όλα τα στοιχεία.

\textbf{Αλγόριθμος.}
\begin{codebox}
\begin{verbatim}
Διαδικασία συγχώνευση(a, b):
  i <- 1
  j <- 1
  Για k = 1 έως m + n:
    Εάν a_i < b_j:
      c_k <- a_i
      i <- i + 1
    Αλλιώς:
      c_k <- b_j
      j <- j + 1
\end{verbatim}
\end{codebox}
(Με κατάλληλους ελέγχους ορίων όταν εξαντληθεί ένας από τους δύο πίνακες ---
αντίγραψε ό,τι μένει από τον άλλον.)

\textbf{Διαίσθηση.} Δύο δείκτες ($i, j$) σαρώνουν τους δύο πίνακες ξεκινώντας από
την αρχή. Σε κάθε βήμα συγκρίνουμε τα κορυφαία στοιχεία και τραβάμε το μικρότερο
στο αποτέλεσμα. Σαν δύο σειρές στο supermarket να γίνονται μία.

\textbf{Πολυπλοκότητα.} Κάθε βήμα του βρόχου είναι σταθερού κόστους (μία σύγκριση
$+$ μία ανάθεση $+$ μία αύξηση δείκτη). Ο βρόχος εκτελείται $m+n$ φορές $\to
O(m+n)$, δηλαδή γραμμικός στο σύνολο των στοιχείων.

\textbf{Χρόνος μνήμης.} Χρειάζεται ένας επιπλέον πίνακας μεγέθους $m+n$ --- η
συγχώνευση \textbf{δεν} γίνεται in-place στην απλή της εκδοχή.
\end{examplebox}

\section{Ο αλγόριθμος mergesort}

Με τη συγχώνευση στο χέρι, ο πλήρης αλγόριθμος είναι σχεδόν τετριμμένος:

\begin{examplebox}
\textbf{Mergesort (αναδρομική).}
\begin{codebox}
\begin{verbatim}
Διαδικασία mergesort(A, p, r):
  Εάν p = r:
    επίστρεψε A[p..r]            // βάση: ένα στοιχείο, ήδη ταξινομημένο
  Εάν p < r:
    q <- floor((p + r) / 2)      // σπάμε στη μέση
    mergesort(A, p, q)           // ταξινομούμε αναδρομικά το αριστερό μισό
    mergesort(A, q + 1, r)       // ταξινομούμε αναδρομικά το δεξί μισό
    επίστρεψε συγχώνευση(A[p..q], A[q+1..r])
\end{verbatim}
\end{codebox}
\textbf{Παρατήρηση.} Καλούμε αρχικά \texttt{mergesort(A, 1, n)}.

\textbf{Γιατί στη μέση;} Θα μπορούσαμε να διαιρέσουμε σε «$n-2$ $+$ $2$» αντί για
«$n/2$ $+$ $n/2$»; \textbf{Όχι} --- η πολυπλοκότητα τότε εκφυλίζεται. Θα το δούμε
στην ανάλυση παρακάτω.
\end{examplebox}

\section{Το δέντρο αναδρομής}

Πριν αποδείξουμε ότι $T(n) = O(n \log n)$, ας δούμε \textbf{οπτικά} πού πάει η
δουλειά. Για $n = 8$, το δέντρο αναδρομής έχει τα εξής επίπεδα:

\begin{center}
\begin{tabular}{@{}l l l l@{}}
\toprule
Επίπεδο & Πλήθος υποπροβλημάτων & Μέγεθος καθενός & Δουλειά επιπέδου \\
\midrule
0 (ρίζα)   & 1 & $n = 8$ & $8$ \\
1          & 2 & $4$     & $2 \times 4 = 8$ \\
2          & 4 & $2$     & $4 \times 2 = 8$ \\
3 (φύλλα)  & 8 & $1$     & $8 \times 1 = 8$ \\
\bottomrule
\end{tabular}
\end{center}

Παρατήρησε το κλειδί: \textbf{κάθε επίπεδο} του δέντρου κάνει \textbf{συνολικά $n$
συγκρίσεις}. Το γιατί;
\begin{itemize}
  \item Στο επίπεδο 1 έχουμε 2 υποπροβλήματα μεγέθους $n/2$. Κάθε ένα κοστίζει
  $n/2$ για συγχώνευση $\to$ σύνολο $n$.
  \item Στο επίπεδο 2 έχουμε 4 υποπροβλήματα μεγέθους $n/4$. Κάθε ένα κοστίζει
  $n/4$ $\to$ σύνολο $n$.
  \item Σε γενικό επίπεδο $k$ έχουμε $2^k$ υποπροβλήματα μεγέθους $n/2^k$. Σύνολο
  $n$.
\end{itemize}
Και υπάρχουν \textbf{$\log_2 n$ επίπεδα} πριν φτάσουμε σε πίνακες μεγέθους 1. Άρα
$T(n) = n \cdot \log_2 n$. Αυτή είναι η διαίσθηση --- τώρα οι τυπικές αποδείξεις.

\section{Πολυπλοκότητα mergesort --- αναδρομική σχέση}

Ορίζουμε $T(n)$ ως την πολυπλοκότητα του αλγορίθμου για είσοδο μεγέθους $n$:
\[
T(n) = \begin{cases}
0 & \text{εάν } n = 1 \\
T(\lfloor n/2 \rfloor) + T(\lceil n/2 \rceil) + n & \text{διαφορετικά}
\end{cases}
\]
Σε λόγια: για $n = 1$ ο πίνακας έχει ένα στοιχείο, καμία πράξη· για $n > 1$ ---
κόστος για το αριστερό μισό $+$ κόστος για το δεξί μισό $+$ $n$ για τη
συγχώνευση.

\textbf{Στόχος:} να δείξουμε ότι $T(n) = O(n \log n)$. Θα παρουσιάσουμε
\textbf{τέσσερις διαφορετικές αποδείξεις} για να εξοικειωθούμε με διάφορες
τεχνικές χειρισμού αναδρομικών σχέσεων. Στην εξέταση, ξέρε \textbf{τουλάχιστον
δύο} από αυτές.

\section{Απόδειξη \#1 --- απλοποίηση της σχέσης (όταν $n = 2^k$)}

Υποθέτουμε ότι το $n$ είναι \textbf{δύναμη του 2}, οπότε δεν χρειάζονται
πατώματα/οροφές. Η σχέση γίνεται $T(n) = 2 T(n/2) + n$.

\textbf{Τέχνασμα.} Διαιρούμε και τα δύο μέλη με $n$ για να πάρουμε ωραίο
τηλεσκοπικό άθροισμα:
\[
\frac{T(n)}{n} = \frac{2 T(n/2)}{n} + 1 = \frac{T(n/2)}{n/2} + 1
\]
Δηλαδή, αν θέσουμε $U(n) = T(n)/n$, τότε $U(n) = U(n/2) + 1$. Εφαρμόζοντας την
ίδια ταυτότητα διαδοχικά:
\[
\begin{aligned}
U(n) &= U(n/2) + 1 \\
     &= U(n/4) + 1 + 1 \\
     &= U(n/8) + 1 + 1 + 1 \\
     &\quad \vdots \\
     &= U(n/n) + \underbrace{1 + 1 + \dots + 1}_{\log_2 n \text{ φορές}}
\end{aligned}
\]
Όταν $n/n = 1$ έχουμε $U(1) = T(1)/1 = 0$, οπότε:
\[
U(n) = \log_2 n \quad \Rightarrow \quad T(n) = n \log_2 n
\]

\begin{keypoint}
\textbf{Αυτή η τεχνική (διαίρεση με $n$ και τηλεσκοπικό άθροισμα)} δουλεύει όποτε
η σχέση είναι $T(n) = 2 T(n/2) + n$. Είναι από τις πιο κομψές αποδείξεις και
αξίζει να την ξέρεις απέξω.
\end{keypoint}

\section{Απόδειξη \#2 --- επαγωγή στον διπλασιασμό}

Με επαγωγή στο $n$ (πάλι υποθέτοντας $n = 2^k$).

\textbf{Βάση.} $T(1) = 0 = 1 \cdot \log 1$. $\checkmark$

\textbf{Επαγωγική υπόθεση.} Έστω $T(n) = n \log n$.

\textbf{Επαγωγικό βήμα.} Θα δείξουμε ότι $T(2n) = 2n \cdot \log(2n)$:
\[
\begin{aligned}
T(2n) &= 2 T(n) + 2n = 2 n \log n + 2n = 2 n (\log n + 1) \\
      &= 2 n (\log n + \log 2) = 2 n \log(2n)
\end{aligned}
\]
όπου χρησιμοποιήσαμε $\log 2 = 1$ (με $\log = \log_2$) και $\log a + \log b =
\log(ab)$. $\blacksquare$

\section{Απόδειξη \#3 --- πολλαπλασιασμός γραμμών (τηλεσκοπικό)}

Ένας ωραίος γενικότερος τρόπος. Για \textbf{οποιαδήποτε} σταθερά $c > 0$, αν
$T(n) \le 2 T(n/2) + c \cdot n$ τότε $T(n) \le c n \log n$.

\textbf{Πώς.} Γράφουμε τη σχέση επανειλημμένα, \textbf{πολλαπλασιάζοντας κάθε
γραμμή με $2^{k-1}$} για να εμφανιστεί τηλεσκόπιο:
\[
\begin{aligned}
T(n)              &\le 2 T(n/2) + cn \\
2 T(n/2)          &\le 4 T(n/4) + 2 c (n/2) = 4 T(n/4) + cn \\
4 T(n/4)          &\le 8 T(n/8) + 4 c (n/4) = 8 T(n/8) + cn \\
                  &\quad \vdots \\
2^{k-1} T(n/2^{k-1}) &\le 2^k T(n/2^k) + cn
\end{aligned}
\]
\textbf{Προσθέτοντας} όλες τις γραμμές, οι μεσαίοι όροι \textbf{τηλεσκοπούνται}:
\[
T(n) \le 2^k T(n/2^k) + cn \cdot k
\]
Για $n = 2^k$ (οπότε $k = \log_2 n$) και $T(1) = 0$:
\[
T(n) \le 0 + cn \log n = cn \log n
\]
Άρα $T(n) = O(n \log n)$. $\blacksquare$

\section{Απόδειξη \#4 --- γενική περίπτωση (ισχυρή επαγωγή, με οροφές)}

Οι προηγούμενες αποδείξεις υπέθεσαν $n = 2^k$. Τι γίνεται για γενικό $n$ (όχι
δύναμη του 2);

\begin{theorem}
Αν $T(n) = 0$ για $n = 1$ και $T(n) = T(\lfloor n/2 \rfloor) + T(\lceil n/2
\rceil) + n$ για $n > 1$, τότε $T(n) \le n \lceil \log n \rceil$.
\end{theorem}

\textbf{Απόδειξη με ισχυρή επαγωγή στο $n$.}

\textbf{Βάση} ($n = 1$): $T(1) = 0 \le 1 \cdot \lceil \log 1 \rceil = 0$.
$\checkmark$

\textbf{Επαγωγική υπόθεση.} Έστω ότι η σχέση ισχύει για κάθε $k = 1, 2, \dots,
n-1$.

\textbf{Επαγωγικό βήμα.} Δείχνουμε ότι $T(n) \le n \lceil \log n \rceil$. Για
ευκολία θέτουμε $n_1 = \lfloor n/2 \rfloor$ και $n_2 = \lceil n/2 \rceil$. Είναι
$n_1 + n_2 = n$ και $n_1 \le n_2$.
\[
\begin{aligned}
T(n) &= T(n_1) + T(n_2) + n \\
     &\le n_1 \lceil \log n_1 \rceil + n_2 \lceil \log n_2 \rceil + n
        && \text{(επαγ.\ υπόθεση)} \\
     &\le n_1 \lceil \log n_2 \rceil + n_2 \lceil \log n_2 \rceil + n
        && (n_1 \le n_2) \\
     &= n \lceil \log n_2 \rceil + n \\
     &\le n (\lceil \log n \rceil - 1) + n && (\star) \\
     &= n \lceil \log n \rceil
\end{aligned}
\]
\textbf{Δικαιολόγηση του $(\star)$.} Παρατηρούμε:
\[
n_2 = \lceil n/2 \rceil \le \frac{2^{\lceil \log n \rceil}}{2} = 2^{\lceil \log n
\rceil - 1}
\]
(γιατί $n \le 2^{\lceil \log n \rceil}$ από τον ορισμό της οροφής). Άρα
$\log n_2 \le \lceil \log n \rceil - 1$, και επειδή το αριστερό μέλος γίνεται
ακέραιος όταν παίρνουμε οροφή, $\lceil \log n_2 \rceil \le \lceil \log n \rceil -
1$. $\blacksquare$

\begin{warningbox}
\textbf{Γιατί χρειαζόμαστε αυτή την «βαρετή» απόδειξη;} Γιατί οι προηγούμενες
υπέθεταν $n = 2^k$. Στην πραγματικότητα το $n$ μπορεί να είναι \textbf{οτιδήποτε}
(π.χ.\ 100, 1000, $13.547$). Η απόδειξη \#4 σου δίνει εγγύηση \textbf{για κάθε
$n$}. Στις εξετάσεις, αν δεν λέει διαφορετικά, χρησιμοποίησε την απλούστερη
απόδειξη ($n = 2^k$) --- είναι αποδεκτή για κάτω/άνω φράγμα.
\end{warningbox}

\section{Γιατί στη μέση και όχι «$n-2$ και $2$»;}

Πρέπει να εξηγήσουμε γιατί διαιρούμε σε ίσα μέρη. Αν διαιρούσαμε σε «μεγάλο $+$
μικρό» --- π.χ.\ $T(n) = T(n-2) + T(2) + n$ --- τότε $T(n) = T(n-2) +
\Theta(n)$. Λύνοντας:
\[
T(n) = \Theta(n) + \Theta(n-2) + \Theta(n-4) + \dots = \Theta(n^2)
\]
\textbf{Δηλαδή χάσαμε όλο το όφελος του D\&C} --- γυρίσαμε στο $\Theta(n^2)$ της
bubble sort. Το «\textbf{στη μέση}» είναι κρίσιμο: μόνο τότε το βάθος του δέντρου
είναι $\log n$ και όχι $n$.

\section{Γενίκευση: άνισος αλλά αναλογικός διαχωρισμός}

\textbf{Σπόιλερ:} το \textbf{βάθος} του δέντρου είναι $\log n$ ακόμα και αν
χωρίσουμε σε \textbf{οποιαδήποτε σταθερή αναλογία} --- όχι μόνο 50/50.

\begin{examplebox}
\textbf{Παράδειγμα.} $T(n) = T(n/3) + T(2n/3) + cn$.

\textbf{Δέντρο αναδρομής.} Το αριστερό υποδέντρο σπάει $1/3$ κάθε φορά, το δεξί
$2/3$. Το βαθύτερο μονοπάτι (πάντα δεξιά) έχει βάθος $\log_{3/2} n$ --- ακόμα
\textbf{λογαριθμικό}. Κάθε επίπεδο κοστίζει $cn$ ακριβώς όπως πριν (το σύνολο
μεγεθών στο επίπεδο είναι πάντα $n$). Άρα συνολικά $T(n) = O(n \log n)$, ακριβώς
όπως η ισόρροπη mergesort.
\end{examplebox}

\begin{keypoint}
\textbf{Γενικό αξίωμα.} Όσο ο διαχωρισμός είναι \textbf{αναλογικός} (κάθε
υποπρόβλημα είναι \textbf{σταθερό κλάσμα} του γονικού), το βάθος είναι $O(\log n)$
και η συνολική δουλειά $O(n \log n)$ --- όταν κάθε επίπεδο κοστίζει $O(n)$.
\end{keypoint}

\section{Άλλο D\&C: ο πύργος του Hanoi}

D\&C δεν είναι μόνο για ταξινόμηση. Ένα κλασικό παράδειγμα:

\begin{examplebox}
\textbf{Πρόβλημα.} Έχουμε 3 στύλους (1, 2, 3) και $n$ δίσκους τοποθετημένους στον
στύλο 1 σε \textbf{κωνική μορφή} (μεγάλος κάτω, μικρός πάνω). Πρέπει να
μεταφέρουμε \textbf{όλους τους δίσκους στον στύλο 3}, διατηρώντας την κωνική
μορφή.

\textbf{Κανόνες.}
\begin{enumerate}
  \item Σε κάθε γύρο μετακινούμε \textbf{μόνο έναν δίσκο}.
  \item Σε κάθε κίνηση αφαιρούμε τον δίσκο που βρίσκεται \textbf{στην κορυφή}
  κάποιου στύλου και τον τοποθετούμε στην κορυφή \textbf{κάποιου άλλου} στύλου.
  \item \textbf{Δεν επιτρέπεται} ένας δίσκος να τοποθετηθεί πάνω σε άλλον με
  \textbf{μικρότερη διάμετρο}.
\end{enumerate}
\end{examplebox}

\subsection{D\&C λύση}

\begin{enumerate}
  \item \textbf{Βάση} ($n = 1$): μία κίνηση --- μεταφορά του δίσκου από Στύλο 1
  $\to$ Στύλο 3.
  \item \textbf{Επαγωγή.} Έστω ότι ξέρουμε να λύνουμε για $n-1$ δίσκους. Για $n$
  δίσκους: μετάφερε τους $n-1$ πάνω δίσκους από Στύλο 1 $\to$ Στύλο 2 (αναδρομική
  κλήση, με Στύλο 3 ως βοηθητικό)· μετάφερε τον $n$-οστό (μεγαλύτερο) δίσκο από
  Στύλο 1 $\to$ Στύλο 3 (μία κίνηση)· μετάφερε τους $n-1$ δίσκους από Στύλο 2
  $\to$ Στύλο 3 (πάλι αναδρομική κλήση).
\end{enumerate}

\textbf{Πολυπλοκότητα.} Έστω $T(n)$ ο αριθμός κινήσεων:
\[
T(n) = T(n-1) + 1 + T(n-1) = 2 T(n-1) + 1
\]
με $T(1) = 1$. Λύνοντας: $T(n) = 2^n - 1$. (Επαλήθευση: $T(1) = 1$,
$T(2) = 2 \cdot 1 + 1 = 3$, $T(3) = 2 \cdot 3 + 1 = 7 = 2^3 - 1$. $\checkmark$)

\begin{warningbox}
\textbf{Η πικρή διαφορά με τη mergesort.} Η mergesort έχει $T(n) = O(n \log n)$
--- \textbf{καλή} πολυπλοκότητα. Ο πύργος Hanoi έχει $T(n) = 2^n - 1$ ---
\textbf{κακή}, εκθετική. Παρά το ίδιο σχήμα (D\&C), η αναλογία διαχωρισμού δεν
είναι ίδια: στο Hanoi κάθε επίπεδο \textbf{διπλασιάζεται}, ενώ στη mergesort το
συνολικό κόστος ανά επίπεδο μένει $n$.

\textbf{Ηθικό δίδαγμα.} Το σχήμα D\&C \textit{από μόνο του} δεν εγγυάται γρήγορο
αλγόριθμο. Πρέπει επίσης το συνδυαστικό βήμα να είναι «φτηνό» σε σχέση με το
μέγεθος των υποπροβλημάτων.
\end{warningbox}

\section{Το Master Theorem}

Είδαμε ότι αναδρομικές σχέσεις της μορφής $T(n) = a T(n/b) + f(n)$ προκύπτουν
συνεχώς. Αντί να αποδεικνύουμε κάθε φορά από την αρχή, υπάρχει ένα \textbf{γενικό
θεώρημα} που τις λύνει.

\begin{keypoint}
\textbf{Master Theorem (Θεώρημα Κυρίαρχου Όρου).} Έστω $a \ge 1$, $b > 1$,
$d \ge 0$ σταθερές και $T(n) = a \cdot T(n/b) + O(n^d)$. Τότε:
\[
T(n) = \begin{cases}
O(n^d)          & \text{εάν } d > \log_b a \\
O(n^d \log n)   & \text{εάν } d = \log_b a \\
O(n^{\log_b a}) & \text{εάν } d < \log_b a
\end{cases}
\]
\end{keypoint}

Σε λόγια: συγκρίνεις τον εκθέτη $d$ του «συνδυαστικού βήματος» με τον $\log_b a$
(τον εκθέτη του δέντρου αναδρομής). Όποιος είναι μεγαλύτερος \textbf{κυριαρχεί}.

\subsection{Διαισθητική απόδειξη --- δέντρο αναδρομής}

Στο $k$-οστό επίπεδο του δέντρου αναδρομής υπάρχουν $a^k$ υποπροβλήματα, καθένα
μεγέθους $n/b^k$, καθένα με $O((n/b^k)^d)$ συνδυαστική δουλειά. Άρα η συνολική
δουλειά στο επίπεδο $k$:
\[
a^k \cdot O\!\left(\left(\frac{n}{b^k}\right)^d\right) = O(n^d) \cdot
\left(\frac{a}{b^d}\right)^k
\]
Άθροισμα από επίπεδο 0 ως $\log_b n$:
\[
T(n) = O(n^d) \cdot \sum_{k=0}^{\log_b n} \left(\frac{a}{b^d}\right)^k
\]
Το τελικό αποτέλεσμα εξαρτάται από τον λόγο $r = a / b^d$:
\begin{itemize}
  \item $r < 1$ ($d > \log_b a$): η δουλειά «μειώνεται» από επίπεδο σε επίπεδο.
  \textbf{Το πρώτο επίπεδο (ρίζα) κυριαρχεί.} Η γεωμετρική σειρά συγκλίνει $\to
  O(n^d)$.
  \item $r = 1$ ($d = \log_b a$): η δουλειά είναι \textbf{ίση σε κάθε επίπεδο}
  (όπως στη mergesort). $\log n$ επίπεδα $\times\ n^d \to O(n^d \log n)$.
  \item $r > 1$ ($d < \log_b a$): η δουλειά «αυξάνεται» από επίπεδο σε επίπεδο.
  \textbf{Τα φύλλα κυριαρχούν.} Το άθροισμα κυριαρχείται από τον τελευταίο όρο
  $\to O(n^{\log_b a})$.
\end{itemize}

\subsection{Πλήρης υπολογισμός}

Το βασικό άθροισμα:
\[
\sum_{k=0}^{\log_b n} r^k = \begin{cases}
O(1)            & r < 1 \\
\log_b n + 1    & r = 1 \\
O(r^{\log_b n}) & r > 1
\end{cases}
\]
Για $r > 1$: $r^{\log_b n} = (a/b^d)^{\log_b n} = a^{\log_b n} / b^{d \log_b n} =
n^{\log_b a} / n^d$. Άρα συνολικά:
\[
T(n) = O(n^d) \cdot \sum_{k=0}^{\log_b n} r^k = \begin{cases}
O(n^d) \cdot O(1) = O(n^d) & r < 1 \\
O(n^d)(\log_b n + 1) = O(n^d \log n) & r = 1 \\
O(n^d) \cdot O(n^{\log_b a}/n^d) = O(n^{\log_b a}) & r > 1
\end{cases}
\]
Που είναι ακριβώς το \textbf{Master Theorem}.

\section{Εφαρμογές του Master Theorem}

\begin{examplebox}
\textbf{Mergesort.} $T(n) = 2 T(n/2) + n$. Άρα $a = 2,\ b = 2,\ d = 1$. Είναι
$\log_b a = \log_2 2 = 1 = d$ --- \textbf{Case 2}:
\[
T(n) = O(n^1 \log n) = O(n \log n) \quad \checkmark
\]
Μας δίνει την ίδια απάντηση όπως οι αποδείξεις \#1--\#4 παραπάνω, σε \textbf{μία
γραμμή}.
\end{examplebox}

\begin{examplebox}
\textbf{Δυαδική αναζήτηση.} $T(n) = T(n/2) + 1$. Άρα $a = 1,\ b = 2,\ d = 0$.
Είναι $\log_b a = \log_2 1 = 0 = d$ --- \textbf{Case 2}:
\[
T(n) = O(n^0 \log n) = O(\log n)
\]
\end{examplebox}

\begin{examplebox}
\textbf{Karatsuba} (θα δούμε στο επόμενο μάθημα). $T(n) = 3 T(n/2) + n$. Άρα
$a = 3,\ b = 2,\ d = 1$. Είναι $\log_b a = \log_2 3 \approx 1.585 > 1 = d$ ---
\textbf{Case 3}:
\[
T(n) = O(n^{\log_2 3}) = O(n^{1.585})
\]
Που είναι \textbf{καλύτερο} από το naive $\Theta(n^2)$ του πολλαπλασιασμού!
\end{examplebox}

\begin{examplebox}
\textbf{Επιλογή --- Quickselect.} $T(n) = T(n/2) + n$. Άρα $a = 1,\ b = 2,\
d = 1$. Είναι $\log_b a = 0 < 1 = d$ --- \textbf{Case 1}:
\[
T(n) = O(n^1) = O(n)
\]
Γραμμική επιλογή του $k$-στού στοιχείου --- θα τη δούμε αργότερα.
\end{examplebox}

\begin{keypoint}
\textbf{Πρακτικός κανόνας για εξεταστική.} Όταν βλέπεις $T(n) = a T(n/b) + n^d$:
(1) υπολόγισε $\log_b a$· (2) σύγκρινέ το με $d$· (3) διάλεξε case και γράψε το
αποτέλεσμα. Τρεις γραμμές --- πιο γρήγορο από κάθε άλλη μέθοδο όταν εφαρμόζεται.
\end{keypoint}

\section{Περιορισμοί του Master Theorem}

Το Master Theorem \textbf{δεν} καλύπτει κάθε αναδρομική σχέση:
\begin{enumerate}
  \item \textbf{Άνισος διαχωρισμός:} $T(n) = T(n/3) + T(2n/3) + n$ $\to$ λύνεται
  με δέντρο αναδρομής, όχι Master.
  \item \textbf{Μη πολυωνυμικό $f(n)$:} $T(n) = 2T(n/2) + n \log n$ $\to$ δεν
  είναι της μορφής $n^d$. Χρειάζεται γενίκευση (Master Theorem υπό μορφή Cormen).
  \item \textbf{$a < 1$ ή $b \le 1$:} δεν έχει νόημα στο μοντέλο.
  \item \textbf{Μη ακέραιο μέγεθος υποπροβλημάτων:} πατώματα/οροφές. Η
  ασυμπτωτική απάντηση μένει ίδια --- αλλά η ακριβής σταθερά αλλάζει.
\end{enumerate}
Στις περιπτώσεις όπου το Master Theorem δεν εφαρμόζεται, \textbf{επιστρέφουμε στις
τεχνικές \#1--\#4} που είδαμε για τη mergesort.

\section{Είναι η mergesort ασυμπτωτικά βέλτιστη;}

Στο δεύτερο μάθημα ορίσαμε «ασυμπτωτικά βέλτιστο» αλγόριθμο. Για την ταξινόμηση
\textbf{με συγκρίσεις}, υπάρχει το γνωστό αποτέλεσμα:

\begin{quoteblock}
\textbf{Θεώρημα (κάτω φράγμα).} Κάθε αλγόριθμος ταξινόμησης που χρησιμοποιεί
\textbf{μόνο συγκρίσεις} στοιχείων απαιτεί $\Omega(n \log n)$ συγκρίσεις στη
χείριστη περίπτωση.
\end{quoteblock}

Η απόδειξη είναι μέσω \textbf{δέντρου απόφασης}: υπάρχουν $n!$ δυνατές
μεταθέσεις, άρα το δέντρο έχει $\ge n!$ φύλλα, άρα ύψος $\ge \log_2(n!) =
\Theta(n \log n)$.

\textbf{Συμπέρασμα:} mergesort $O(n \log n)$ άνω· κάθε ταξινόμηση με συγκρίσεις
$\Omega(n \log n)$ κάτω. Άρα η mergesort είναι \textbf{ασυμπτωτικά βέλτιστη} στην
οικογένεια αλγορίθμων ταξινόμησης με συγκρίσεις.

\begin{intuition}
Η μόνη μέθοδος να ξεπεράσεις το $\Omega(n \log n)$ είναι να \textbf{μην}
χρησιμοποιήσεις μόνο συγκρίσεις --- π.χ.\ counting sort / radix sort, που
χρησιμοποιούν δομή των δεδομένων (ακέραιοι σε γνωστό εύρος). Αυτοί τρέχουν σε
$O(n)$ αλλά λειτουργούν μόνο σε ειδικές περιπτώσεις. Στο μάθημα μένουμε στην
ταξινόμηση με συγκρίσεις.
\end{intuition}

\begin{recapbox}
\textbf{Τι κρατάμε από αυτή τη διάλεξη}
\begin{enumerate}
  \item \textbf{D\&C ως σχήμα}: διαίρει σε ξένα μικρότερα υποπροβλήματα $\to$
  κυρίευσε αναδρομικά $\to$ συνδύασε. Δικαιολογείται με επαγωγή.
  \item \textbf{Mergesort}: σπάμε στη μέση, αναδρομικά ταξινομούμε, συγχωνεύουμε
  με γραμμικό κόστος $\to O(n \log n)$. \textbf{Τέσσερις διαφορετικές
  αποδείξεις}, κάθε μία χρήσιμη σε διαφορετική εξεταστική.
  \item \textbf{Γιατί στη μέση}: μόνο αναλογικός διαχωρισμός δίνει $\log n$
  βάθος. «$n-2$ $+$ $2$» κάνει τον αλγόριθμο $\Theta(n^2)$.
  \item \textbf{Πύργος Hanoi}: D\&C με \textbf{εκθετική} πολυπλοκότητα
  ($2^n - 1$). Το σχήμα από μόνο του δεν εγγυάται γρήγορο αλγόριθμο.
  \item \textbf{Master Theorem}: για $T(n) = a T(n/b) + O(n^d)$ ---
  $d > \log_b a \to O(n^d)$· $d = \log_b a \to O(n^d \log n)$·
  $d < \log_b a \to O(n^{\log_b a})$.
  \item \textbf{Η mergesort είναι ασυμπτωτικά βέλτιστη} για ταξινόμηση με
  συγκρίσεις (κάτω φράγμα $\Omega(n \log n)$).
\end{enumerate}
\end{recapbox}

\lecturefooter
\end{document}
